\documentclass[10pt,a4paper]{article}
\newcommand{\runningtitle}{Held-out population-flow evaluation in acute human wounds}
\usepackage[margin=15mm,top=16mm,bottom=17mm,headheight=12pt,headsep=5mm,footskip=9mm]{geometry}
\usepackage{fontspec}
\setmainfont{TeX Gyre Termes}
\setsansfont{Arial}
\usepackage{amsmath,amssymb,mathtools,bm}
\usepackage{unicode-math}
\setmathfont{TeX Gyre Termes Math}
\usepackage{graphicx,booktabs,array,tabularx,longtable}
\usepackage[font=small,labelfont=bf,labelsep=period,skip=4pt]{caption}
\usepackage{flafter,placeins,needspace}
\usepackage[sort&compress,numbers]{natbib}
\usepackage[unicode,colorlinks=true,linkcolor=black,citecolor=black,urlcolor=black]{hyperref}
\usepackage{xcolor,microtype,titlesec,enumitem,fancyhdr}
\usepackage{bookmark}
\hypersetup{pdfcreator={XeLaTeX},pdfauthor={Zeyu Fu}}
\graphicspath{{../figures/}}
\setlength{\parindent}{1em}
\setlength{\parskip}{2pt plus .5pt}
\linespread{1.035}
\setlength{\textfloatsep}{9pt plus 2pt minus 2pt}
\setlength{\floatsep}{8pt plus 2pt minus 2pt}
\setlength{\intextsep}{8pt plus 2pt minus 2pt}
\renewcommand{\topfraction}{.94}
\renewcommand{\bottomfraction}{.88}
\renewcommand{\textfraction}{.06}
\renewcommand{\floatpagefraction}{.78}
\setcounter{topnumber}{3}
\setcounter{bottomnumber}{2}
\setcounter{totalnumber}{4}
\setlength{\abovedisplayskip}{6pt plus 2pt minus 2pt}
\setlength{\belowdisplayskip}{6pt plus 2pt minus 2pt}
\setlength{\abovedisplayshortskip}{3pt plus 1pt}
\setlength{\belowdisplayshortskip}{4pt plus 1pt}
\titlespacing*{\section}{0pt}{12pt plus 2pt minus 2pt}{5pt}
\titlespacing*{\subsection}{0pt}{8pt plus 2pt minus 1pt}{3pt}
\titleformat{\section}{\large\bfseries}{\thesection}{.6em}{}
\titleformat{\subsection}{\normalsize\bfseries}{\thesubsection}{.6em}{}
\setlist{nosep,leftmargin=1.7em,topsep=4pt}
\setlength{\bibsep}{1.5pt}
\renewcommand{\bibfont}{\small}
\setlength{\tabcolsep}{4pt}
\setlength{\LTpre}{2pt}
\setlength{\LTpost}{2pt}
\renewcommand{\arraystretch}{1.12}
\setlength{\emergencystretch}{1.2em}
\widowpenalty=10000
\clubpenalty=10000
\flushbottom
\pagestyle{fancy}
\fancyhf{}
\fancyhead[L]{\sffamily\fontsize{8}{10}\selectfont\runningtitle}
\fancyfoot[C]{\small\thepage}
\renewcommand{\headrulewidth}{.25pt}
\newcommand{\papertitle}[1]{\thispagestyle{plain}\begin{center}\vspace*{-5mm}{\fontsize{19}{22}\selectfont\bfseries #1\par}\vspace{2mm}{\normalsize Zeyu Fu\par}\end{center}\vspace{-2mm}}
\newcommand{\paperfigure}[3]{\begin{figure}[!htbp]\centering\includegraphics[width=\linewidth]{#1.pdf}\caption{#3}\label{#2}\end{figure}}
\newcommand{\supplement}{\FloatBarrier\section*{Supplementary methods and numerical results}\addcontentsline{toc}{section}{Supplementary methods and numerical results}\setcounter{section}{0}\setcounter{table}{0}\setcounter{equation}{0}\renewcommand{\thesection}{S\arabic{section}}\renewcommand{\thetable}{S\arabic{table}}\renewcommand{\theequation}{S\arabic{equation}}}
\newcolumntype{L}{>{\raggedright\arraybackslash}X}
\newcommand{\doi}[1]{\href{https://doi.org/#1}{doi: #1}}
\newcommand{\E}{\mathbb{E}}
\newcommand{\R}{\mathbb{R}}
\newcommand{\ind}{\mathbf{1}}

\hypersetup{pdftitle={Population flow reconstruction of an acute human wound interval: held-out evaluation, time coordinates and donor-count sensitivity}}
\begin{document}
\papertitle{Population flow reconstruction of an acute human wound interval: held-out evaluation, time coordinates and donor-count sensitivity}
\section*{Abstract}
\noindent A smooth trajectory fitted to cross-sectional single-cell data does not establish prediction of an individual's wound course. We tested whether population fields improve on no motion when a target time point or donor is excluded from fitting.
Three healthy donors contributed intact skin and acute wounds at days 1, 7 and 30. A frozen expression representation, training-only standardization and donor-paired endpoint sampling supported separate timepoint-holdout and donor-transfer protocols. We compared four time coordinates, shared and donor-conditioned flow matching, three non-neural predictors, two training-donor counts and three training seeds. Three mouse model arms supplied a separate computational audit.
Pooled Wound7 reconstruction reduced energy distance relative to stay-still by +49.9\%. Excluding each donor preserved the result for the retained donors in 3/3 sensitivity runs; genuine held-out-donor transfer separately improved by a mean +75.9\% in 3/3 donors. Linear days failed at the target (−18.2\%), whereas rank, square-root and log-day coordinates improved. All five predictive methods met the specified split-half-reference criterion in 3/3 donor folds. Shared flow had the smallest mean error, 0.395; conditioning improved all six training-donor comparisons without improving the overall held-out mean. Adding a second training donor reduced seed-0 mean error from 0.5496 to 0.3854. Aggregate reductions were 0.156--0.176 across seeds, with donor-block interval $[0.081,0.315]$, but some individual configurations varied by 0.368. Mouse-arm changes ranged from −1.5\% to +50.6\%, with one animal per arm and time point.
These experiments support reconstruction of one acute human interval in a three-donor pilot, conditional on the representation and evaluation protocol. They establish neither DFU prognosis nor a population-level method ranking, identifiable biological trajectories, or a required cohort size.
\par\smallskip\noindent Keywords: conditional flow matching; wound healing; held-out donor; energy distance; time coordinate; donor heterogeneity.

\section{Introduction}
Wound repair involves changing mixtures of inflammatory, stromal and epithelial states over time \citep{wilkinson2020}. In diabetic foot disease, clinical healing and recurrence are outcomes of interest, but they cannot be inferred directly from a transcriptional trajectory \citep{armstrong2017}. Destructive single-cell sampling measures distributions at selected times rather than following the same cells longitudinally. Trajectory inference can organize these states, yet its conclusions depend on sampling design, topology and modeling assumptions \citep{saelens2019,weinreb2018}.

Optimal transport and continuous-flow models offer ways to connect observed distributions \citep{schiebinger2019,tong2020trajectory,lipman2023,tong2023}. Neural ordinary differential equations parameterize a continuous field, and flow matching learns such a field from chosen probability paths \citep{chen2018ode,lipman2024guide}. Related transport approaches also predict cellular responses to perturbations \citep{bunne2023}. Reconstructing a target distribution and identifying an individual biological trajectory remain distinct objectives because multiple dynamics can produce the same marginals \citep{lavenant2021,weinreb2018}.

A second distinction concerns generalization. Predicting a time point missing from every donor removes different information from predicting an unseen donor whose target time is already represented in training. Both tasks depend on the expression coordinates and preprocessing; transferable latent representations are useful, but do not eliminate representation choice or sampling effects \citep{lopez2018,heumos2023}. Biological replication also remains at the donor level even when many cells enter a distributional comparison \citep{zimmerman2021}.

The public acute human wound atlas offers a bounded setting in which to separate these questions: three healthy volunteers contribute skin and wound samples at four times \citep{liu2025_wound,geo_gse241132}. We ask whether a predicted distribution approaches a held-out target more closely than an unchanged source distribution. We compare missing-time and missing-donor protocols, time coordinates, neural and non-neural predictors, training-donor counts and training seeds. Numerical-resolution and evaluation-sampling analyses test additional sources of computational sensitivity, while three mouse arms supply a separate model comparison. These experiments assess distribution reconstruction and its limits; the available cohorts do not contain a diabetic-foot-ulcer time-to-closure target.

\section{Results}
\subsection{Observed fibroblasts define the material available for temporal prediction}

Three healthy volunteers supplied intact skin and acute-wound biopsies at days 1, 7 and 30, yielding 12 specimens and 12,259 author-annotated fibroblasts (Figure~\ref{fig:p2f1}A\label{call:p2f1}). The computational route in A retains the 7,002-gene input dimension despite coverage of only 5,383 genes, then distinguishes the frozen mean-coordinate encoder, training-only scaling, product-coupled field fitting and numerical prediction. Its globally held-out day-7 task does not depict biological cell paths or the separate donor-transfer protocol. A common projection displays these observed cells by collection time (Figure~\ref{fig:p2f1}B) and by donor (Figure~\ref{fig:p2f1}C). Its first two principal components explain 43.5\% and 10.6\% of the variance in the training set that excludes day 7. Thus, nearly half of the training variance remains outside the displayed plane; visual overlap or separation cannot replace evaluation in all 15 coordinates.
\paperfigure{figure1_workflow}{fig:p2f1}{Acute-wound sampling and frozen-coordinate flow. (A) Conceptual repair, not measured histology or closure time, and repeated biopsies of three healthy donors at days 0, 1, 7 and 30. The lower route starts with 7,002-gene proportions: 5,383 panel genes are covered and the remainder zero-filled. A frozen 15-dimensional encoder mean μ is scaled using training cells only. A shared conditional-flow-matching (CFM) field learns from independently sampled product-coupled endpoints; 50-step RK4 predicts the day-7 marginal in scaled μ coordinates, not simplex weights. All day-7 cells are held from fitting and scaling. This is global time holdout, not the separate held-donor task that permits training-donor day-7 cells. (B) A uniform 2,000-cell display subset of 12,259 fibroblasts by time. PCA and scaling exclude day 7; the axes explain 43.5\% and 10.6\% of training variance. (C) The same cells by donor. (D) Twelve measured genes per donor--time biopsy, with log-transformed pseudobulk counts per 10,000 scaled separately by gene; colors do not compare absolute levels between genes. (E) Fibroblast counts per biopsy; cells are not independent donors. (F) Thirty-six seed-0 numerical paths from observed day-1 cells to day-7 rank time over the day-1/day-7 display. Arrows are not tracked cell fates; errors use all 15 coordinates.}


Matched raw-count profiles show the expression of matrix, inflammatory and contractile genes in each donor's four biopsies (Figure~\ref{fig:p2f1}D). The gene-wise color scale compares a gene across samples and does not compare absolute expression between genes. Retained fibroblast counts differ across donor--time combinations (Figure~\ref{fig:p2f1}E), so neither point density nor the number of cells provides independent temporal replication. Projected paths connect observed day-1 starting coordinates to model-predicted coordinates at the prespecified day-7 time (Figure~\ref{fig:p2f1}F). These are numerical field trajectories; the data do not track the corresponding cells through the wound. Sampling counts are documented in Table~\ref{tab:p2biologicalcounts}.

% Allow the next Results subsection to fill the text page while Figure 1
% remains on the immediately following float page (checked by build.py).
\subsection{A held-out acute human time point can be reconstructed under the fixed protocol}

Global timepoint holdout removes day 7 from every donor, whereas donor transfer removes one donor and retains all four times in the others. Both use the frozen representation and the training-only scaling in Equation~\eqref{eq:p2scale}. The product bridge in Equation~\eqref{eq:p2bridge}, field objective in Equation~\eqref{eq:p2cfm} and initial-value prediction in Equation~\eqref{eq:p2ode} specify a reproducible estimator, not an observed cell transition. This distinction determines whether the target time itself has been observed during training; the permitted source, target exclusions and fitting settings are recorded in Table~\ref{tab:p2protocols}.

A fixed-seed refit assessed numerical resolution of the RK4 update in Equation~\eqref{eq:p2rk4} and sensitivity to evaluation sampling. Across 25, 50, 100 and 200 RK4 steps, the target energy distance remained 0.705395 and the improvement remained 49.863\%; the largest discrepancy from the 200-step coordinates was below $6\times10^{-7}$ in root-mean-square units (Table~\ref{tab:p2numerics}). Repeated evaluation at 250, 500 and 1,000 cells per distribution gave median improvements of 48.46\%, 48.39\% and 49.26\%, respectively. Gain exceeded the same-count split-half reference in all 50 draws at each count (Table~\ref{tab:p2sampling}). These results assess one fitted field and the observed cells; the ranges across draws quantify evaluation sensitivity, not additional donor replication or confidence in clinical generalization.

Using the empirical weighting rule in Equation~\eqref{eq:p2weighted} with all cells and the same fixed 50-step prediction, equal cell weights gave a 49.40\% improvement and equal donor mass gave 49.69\%. The three donor-specific improvements under this pooled field were 47.15\%, 46.92\% and 40.07\% (Table~\ref{tab:p2weights}). This sensitivity changes the evaluation measure without refitting the field; it is distinct from the donor-transfer experiment and from the pooled capped-distance estimate.

\FloatBarrier
The result is deliberately bounded. The donors are healthy volunteers, the wounds are acute experimental wounds and the target is one observed time point. The analysis does not contain a DFU patient trajectory or a time-to-closure endpoint. ``Transfer'' here means transfer across the three specified acute-wound donors under a held-out protocol; it does not mean clinical generalization.
Using the empirical version (Equation~\eqref{eq:p2edv}) of the energy discrepancy defined in Equation~\eqref{eq:p2ed}, the finalized held-out Wound7 analysis improved on the stay-still baseline by +49.9\% under Equation~\eqref{eq:p2improvement} (Figure~\ref{fig:p2f2}A\label{call:p2f2}). Dropping one donor at a time preserved improvement in 3/3 retained-donor sensitivity runs; these runs predict the retained donors, not the omitted donor. The held-out donor transfer analysis also passed in 3/3 folds, with a mean improvement of +75.9\% (Figure~\ref{fig:p2f2}B). In the geometry audit, the median same-donor temporal distance was 1.967 and the median same-time donor distance was 0.233, giving a time-to-donor ratio of 8.46 (descriptive 95\% resampling range 6.02--13.11; Figure~\ref{fig:p2f2}C). The mean displacement cosine was +0.948, with a descriptive resampling range of 0.943--0.954. Shared donors make these pair distances dependent, so the ranges do not have established donor-population coverage. Improvement over unchanged source alone does not establish incremental value over other permitted observations. The probability-conservation identity in Equation~\eqref{eq:p2continuity} concerns the constructed training paths, not identification of a biological trajectory. The observed reconstruction gains remain conditional on this representation, interval and three donors (Tables~\ref{tab:p2folds} and~\ref{tab:p2geometrytable}).

\paperfigure{figure2_wound7_transfer}{fig:p2f2}{Global Wound7 reconstruction, donor transfer and donor--time geometry. (A) Bars show predicted and stay-still energy distances to pooled Wound7 after excluding that time globally; relative improvement is 100 times (stay-still minus predicted distance) divided by stay-still distance, giving 49.9\%. (B) Each bar is a genuinely held-out donor, with all four times retained in the other two donors; the mean relative improvement is 75.9\%. Bars are point estimates without confidence intervals. (C) Points show 12 between-donor same-time distances and 18 within-donor between-time distances. Boxes span the interquartile range, centre lines are medians, and whiskers extend to extreme points within 1.5 interquartile ranges. The temporal/donor median ratio is 8.46 (10,000-draw descriptive pair-distance bootstrap 95\% interval 6.02--13.11). All distances use training-standardized latent coordinates. Shared donors make these pairs dependent; three healthy donors do not provide DFU prediction validation.}


\subsection{Simple interpolation is competitive when an entire time point is missing}

A retrospective extension tested whether improving over unchanged source requires a neural field. The additional comparisons hold out either day 7 or day 1 globally, keeping the frozen encoder, rank clock, training-only scaling and optimization budget fixed. For each task, the observed time immediately before the target supplies the source and the next observed time supplies an anchor. The source and anchor belong to the same donor; the held-time cells enter only the evaluation. Error is calculated using every available cell in all 15 coordinates within each donor, then averaged with equal donor weight. The donor-averaged error in Equation~\eqref{eq:p2macroerror} differs from the earlier pooled and capped estimands. The two endpoint controls in Equation~\eqref{eq:p2interpolation} therefore address model complexity under this separate, explicit scoring rule.

For withheld day 7, unchanged source had mean donor error 1.4725. Across three seeds, shared flow reduced it to 0.8110, while within-donor centroid translation reduced it further to 0.6674. An independent-endpoint linear bridge gave 0.9284 (Figure~\ref{fig:p2f3}A\label{call:p2f3}). The corresponding improvements were 44.9\%, 54.7\% and 36.9\%. Withholding day 1 instead trained on intact skin, day 7 and day 30 and predicted from intact skin. The unchanged-source error was 1.3460; shared flow, centroid translation and the independent bridge gave 1.0571, 0.9293 and 1.0959, corresponding to improvements of 21.5\%, 31.0\% and 18.6\% (Figure~\ref{fig:p2f3}B). The complete seed-level errors appear in Table~\ref{tab:p2interpolation}.

\paperfigure{figure3_interpolation_controls}{fig:p2f3}{Simple baselines for globally missing times. All methods use the same frozen representation and training exclusions; error averages full-cell empirical energy distances within the three donors. (A) Withheld day 7, predicted from day 1 with day 30 as the observed anchor. (B) Withheld day 1, predicted from intact skin with day 7 as anchor. Points show the three-seed mean error; horizontal ranges show seed minima and maxima, not confidence intervals. Both tasks use the fixed rank-time interpolation fraction one half. (C) Donor-specific seed-averaged field error minus centroid-translation error; positive values favor translation. The later anchor belongs to the same donor, so this is retrospective interpolation rather than source-only donor transfer. (D) Each field seed's improvement over the corresponding unchanged source. Seeds and repeated distance evaluations provide no additional biological replication. (E) Seed- and donor-averaged energy distance after separately centering each distribution; this diagnoses shape without decomposing total error. (F) Mean and full range of flow-minus-comparator errors across 50 paired draws of 200 source and 200 target cells per donor; positive values favor the comparator. Location--scale uses only endpoint means and coordinate standard deviations. Matched draw ranges are not confidence intervals.}


The centroid translation had lower seed-averaged error in all three donors for withheld day 7 and in two of three for withheld day 1 (Figure~\ref{fig:p2f3}C). The flow error ranged from 0.8044 to 0.8237 across day-7 fits, and from 0.9484 to 1.1231 across day-1 fits (Table~\ref{tab:p2interpolation}). Every fit improved over its unchanged source; Figure~\ref{fig:p2f3}D displays these seed-specific relative improvements. These ranges characterize training variation in the existing cohort; they do not provide six independent biological replications. All donor-specific comparisons are retained in Table~\ref{tab:p2interpolationdonors}.

The simple predictor uses an observed later biopsy from the donor being reconstructed, which is permitted by the global timepoint-holdout task. It would not be available in prospective source-only prediction for a new donor. Conversely, the donor-transfer benchmark observes the target time in the training donors and asks a different question. Its method comparisons cannot be generalized to a globally unobserved time. The added comparison supports temporal information in the measured distributions, while providing no evidence that a neural field is necessary for the two interpolation tasks tested here.

The location--scale predictor in Equation~\eqref{eq:p2locationscale}, which interpolates endpoint standard deviations as well as means, gave errors of 0.6982 for day 7 and 0.9034 for day 1, also lower than the shared-flow errors in both tasks (Figure~\ref{fig:p2f3}A--B). Mean translation was the better of these two deterministic controls for day 7, whereas location--scale interpolation was better for day 1. This ordering is reported without selecting a new predictor for either target.

The distribution-shape diagnostics explain why a single aggregate error does not identify a uniform advantage. For day 7, flow had slightly lower centered energy distance than translation (0.1441 versus 0.1495), despite higher total error. For day 1, flow had higher centered error (0.1220 versus 0.0647), while the independent bridge also had higher centered error (0.1445; Figure~\ref{fig:p2f3}E). The independent bridge's predicted-to-target variance ratios were 0.319 and 0.579 at days 7 and 1, indicating substantial contraction in this comparison (Table~\ref{tab:p2shapediagnostics}). These are separate geometry summaries, not a causal or additive error decomposition.

Matched 200-cell evaluation preserved the lower error of both deterministic endpoint controls in all 50 draws for each held time. Flow-minus-translation error ranged from 0.1132 to 0.1716 for day 7 and 0.0977 to 0.1545 for day 1; flow-minus-location--scale ranges were 0.0847--0.1405 and 0.1194--0.1763. Flow remained better than the independent bridge in every draw (Figure~\ref{fig:p2f3}F; Table~\ref{tab:p2matchedsampling}). The full-cell calculations reproduced all previous donor/seed errors to numerical precision. These results establish robustness of the observed ordering to the tested evaluation samples, conditional on three donors and the frozen fits.

\subsection{Time parameterization changes the held-out result}

The historical single-seed, nearest-grid comparison gave linear days an improvement of −18.2\% relative to unchanged source (Table~\ref{tab:p2axes}). The revised paired audit instead evaluates 12 fields at the exact target time with identical random streams and evaluation rows within each seed. The day-based transforms in Equation~\eqref{eq:p2times} define separate fitted bridges, rather than transforming one fixed field by the chain rule in Equation~\eqref{eq:p2clock}. Mean donor energy distances were 0.8104, 2.1542, 0.8746 and 0.7725 for rank, linear days, square-root days and log days, respectively, versus an unchanged-source distance of 1.4746. The corresponding centroid translations gave 0.6667, 0.9775, 0.7476 and 0.6654, lower than the shared field under every clock on the donor-and-seed average (Table~\ref{tab:p2pairedclock}; Figure~\ref{fig:p2f4}\label{call:p2f4}). These equal-donor, exact-time values have a different estimand and sampling cap from the historical pooled scan. Changing the coordinate changes bridges and optimization targets, so the comparison does not isolate interval compression or select a universal biological clock.

\paperfigure{figure4_time_axes}{fig:p2f4}{Paired time-coordinate sensitivity under global day-7 holdout in three healthy donors. Twelve fields combine four fixed clocks with seeds 0, 1 and 2, each using 8,000 training steps and batch size 256. Network initialization, endpoint draws and interpolation random draws are paired across clocks within seed. All day-7 cells are excluded from fitting and scaling. Energy distance is evaluated at the exact transformed day-7 time after 50 RK4 steps, using the same source and target rows across clocks and seeds and a 1,200-cell cap per donor. Each point is the equal mean of three donor-specific distances; shapes identify training seeds. Thin lines join the same seed, and thick lines show three-seed means. Colors distinguish shared flow, within-donor centroid translation and unchanged source. The centroid control uses the permitted day-1 and day-30 biopsies, not day-7 expression. Constant baseline points overlap. Connecting categorical clock choices does not depict a biological trajectory, and seed spread is not a biological confidence interval.}




\subsection{The training-target distribution challenges the held-donor method ranking}

The historical held-donor benchmark gave shared flow matching a mean energy distance of 0.3947 (approximately 0.395) and gains above the split-half reference in 3/3 folds. Mean displacement, nearest-donor displacement, entropic transport and the conditioned field also met that descriptive criterion. An added predictor instead uses only the training donors' Wound7 distributions, assigning equal donor mass and matching the held-source particle count. Its mean error was 0.1818, with a full range of 0.1479--0.2150 across 50 prediction-sampling draws (Figure~\ref{fig:p2f5}A\label{call:p2f5}). The same historical target indices and scaling were used, so improvement over unchanged source does not establish incremental value from the excluded donor's source expression (Tables~\ref{tab:p2methods}, \ref{tab:p2methodsbydonor} and~\ref{tab:p2targetmarginal}).

The donor-specific results prevent a universal ranking: target-marginal mean errors were 0.1569, 0.1416 and 0.2469 for PWH26, PWH27 and PWH28, whereas shared-field errors were 0.7110, 0.2881 and 0.1850. The field therefore remained better in PWH28, including all 50 matched marginal draws. Those draws reuse the same three donors and quantify prediction-sampling sensitivity rather than population uncertainty.

The context model in Equation~\eqref{eq:p2context} improved all 6/6 training-donor comparisons (Figure~\ref{fig:p2f5}B; Table~\ref{tab:p2insample}) while giving a higher held-out mean than the shared field. Its greater fit to known contexts therefore supplies no evidence of a stable generalization advantage.

\paperfigure{figure5_methods_benchmark}{fig:p2f5}{Method comparison across three held-out donors. (A) Bars are equally donor-weighted mean energy distances for unchanged source, five historical predictors and the added matched-particle target marginal. CFM means conditional flow matching; OT is the entropic transport displacement extension. The target marginal assigns equal mass to the two training donors' Wound7 distributions, uses no held-source expression, and averages 50 draws under the original target samples and scaling. Its mean error is 0.1818 versus 0.3947 for historical shared CFM, but PWH28 favors the field. All historical predictive methods exceed the single split-half gain reference in 3/3 folds; this descriptive criterion is not a significance test or population ranking. (B) Lines connect shared and conditioned errors for six training-donor comparisons (two retained donors in each of three fits). Conditioning improves all six. These comparisons share donors and do not measure unseen-donor transfer. Neither panel shows uncertainty intervals.}


The original neural and transport method comparison used one training seed and successive endpoint random draws. The entropic plan in Equation~\eqref{eq:p2ot} is a separate comparator, not the coupling used to train the flow. The added target-marginal result is a retrospective baseline audit, and the historical fields were not saved as per-cell predictions for new paired target sampling. Neither result selects a model class for DFU patients; independent donors and matched repeated fits are required for that question.

\subsection{Donor-count comparisons require a common evaluation scale}

The initial count curve changed its standardization when the training subset changed, so subtracting its absolute errors mixed information gain with a change of distance geometry. A controlled rerun retained subset-only training normalization but scored inverse-transformed predictions in one common reference per held donor. Across 27 fits, the hierarchical average in Equation~\eqref{eq:p2curve} gave shared-field mean error of 0.5466 with one training donor and 0.3983 with two. Figure~\ref{fig:p2f6}A\label{call:p2f6} shows all three seeds and the seed-0 line (0.5315 to 0.3854). Seed-specific reductions were 0.1462, 0.1558 and 0.1430 (Figure~\ref{fig:p2f6}B). The mean reduction was 0.1483, with all 3/3 donor blocks improving and a descriptive three-block percentile range of {[}0.0732, 0.2727{]} (Table~\ref{tab:p2curve}).

\paperfigure{figure6_donor_curve_stability}{fig:p2f6}{Donor-count comparison in common scoring coordinates. (A) k is the number of training donors. The line joins seed-0 mean held-out errors, 0.5315 at k=1 and 0.3854 at k=2; symbols identify seeds 0, 1 and 2. Each model uses only its selected training donors for fitting and normalization. Predictions return to encoder coordinates before scoring under one reference fitted from both non-held donors, fixed across k and seed and excluded from single-donor fitting. At k=1, two training choices are averaged within held donor before equally weighting the three donors. (B) Points are seed-specific reductions, k=1 minus k=2. The line is the three-seed mean, 0.1483. Shading is a descriptive percentile 95\% range [0.0732, 0.2727] from 4,000 resamples of three donor blocks after seed averaging; it is not an interval for each point. Source and target indices are fixed. The largest configuration-specific seed range is 0.3377. Overlapping training sets limit interval interpretation, and two k values cannot identify a plateau or required cohort size.}


The same common geometry also showed lower error for mean displacement (0.5095 to 0.4186), the full equal-donor target marginal (0.2222 to 0.1676), and its matched-particle version (0.2332 to 0.1818). Unchanged-source error remained 1.4760 at both training sizes, checking that the reference metric did not change with k (Table~\ref{tab:p2seedcurve}). The full target-marginal decrease, however, is exactly the mixture effect in Equation~\eqref{eq:p2mixture}; it cannot independently identify target-relevant information from adding a donor. The neural and mean-displacement comparisons involve different fitted predictors and are not algebraically identical to this target-marginal mixture.

The saved two-donor fields also allow the starting distribution to be changed without changing training, target cells or scoring scale. Averaged across seeds and donors, the actual held-source prediction had error 0.3983. Replacing its input with equal-donor draws of training Wound1 cells reduced mean error to 0.2450; the matched training Wound7 marginal, which does not use any source expression from the held donor, had error 0.1818 (Table~\ref{tab:p2sourceswap}). The donor effects were heterogeneous: real-source versus substituted-source errors were 0.7586 versus 0.2802 for PWH26, 0.2368 versus 0.2732 for PWH27, and 0.1994 versus 0.1816 for PWH28. Thus, PWH27 favored its own starting cells on average, while the other two did not. The substitution is a sensitivity analysis of the same fitted fields, not a comparison of nine independent donors; the 50 cell draws likewise supply no new biological replication. It does not demonstrate a consistent incremental value for the excluded donor's initial expression in this cohort, and it does not imply that such information is universally unhelpful.

The exact target-marginal reduction was 0.054601 across donors and matched one quarter of the training-target mutual discrepancy, with maximum numerical residual below $2.45\times10^{-15}$ (Table~\ref{tab:p2mixtureidentity}). The additional 200-particle source audit gave macro errors of 0.4036 for real sources, 0.3075 when the two individual substituted-source errors were averaged, and 0.2527 for their sampled mixture. PWH28 exposes the distinction: real-source error was 0.2043, lower than either individual substitution (0.2503 and 0.2636), yet the mixed substitution was lower still at 0.1904. PWH27 also favored its real source over both individual alternatives, while PWH26 did not (Table~\ref{tab:p2mixtureaudit}). Thus, the earlier mixed-source comparison combines source replacement with a change in the predicted mixture and cannot isolate individual source-state value. The exact identity does not decompose the gain from a jointly fitted two-donor field. All results remain conditional on the same three donors, fields, target cells and reference geometry.

Within a fixed donor and training subset, the mean training SD was 0.0701, or 3.07 times the mean split-half reference of 0.0228; the largest seed range was 0.3377 (Table~\ref{tab:p2seedconfigs}). The aggregate direction persisted across seeds, while individual fits remained variable. The three-block range is descriptive because the held-donor folds share training donors. Two training sizes cannot determine a plateau or a required cohort size, and seeds, subsets and repeated marginal draws add no independent biological replication.

\subsection{Observable expression errors do not establish a flow advantage over target marginals}

Raw-expression alignment and frozen re-encoding passed for all 12,259 cells, with maximum encoder-mean discrepancy $1.91\times10^{-6}$. On the exact saved targets and all 5,383 common-panel genes, shared-flow total variation was 0.2934 with one training donor and 0.2823 with two, versus 0.3529 for decoded persistence. The decoded mean-displacement control gave 0.2882 and 0.2825; decoded training-target marginals gave 0.2590 and 0.2546. Directly observed training-target marginals were lower still at 0.1279 and 0.1115, without using the excluded donor's source expression (Table~\ref{tab:p2observable}). Flow improved over decoded persistence in all three donor blocks but over decoded target marginals in only one of three, and over observed target marginals in none. At two training donors its total variation was close to mean displacement, while its Jensen--Shannon divergence was higher, so these metrics do not establish a consistent neural advantage. The all-day-7 sensitivity retained the same aggregate interpretation. The observed-target decoder reference had mean total variation 0.2430 against actual expression (Table~\ref{tab:p2decoderreference}); because it uses the target itself, this is neither prediction performance nor an error floor. The results add a mean-expression evaluation within the existing cohort, not independent biological validation or evidence about unassayed genes.

\subsection{Three mouse arms expose model-arm heterogeneity}

The all-arm audit gave different results under the same protocol, using the endpoint-distance diagnostic in Equation~\eqref{eq:p2geometry}. NDB changed by −1.5\% relative to stay-still and did not beat the baseline. PDB improved by +32.1\%, but its POD7 target failed an endpoint-distance ordering check, warning against a simple interpolation interpretation without proving an off-manifold trajectory. GDB improved by +50.6\% with a POD7 target satisfying that ordering check. In aggregate, 2/3 arms beat stay-still, 2/3 met the noise criterion and 2/3 satisfied the ordering check (Figure~\ref{fig:p2f7}A\label{call:p2f7}). Each arm had one animal per time point, and NDB covered 59.7\% of the frozen panel. These arm-level computations are useful for protocol auditing, but they do not establish a mouse biological trajectory, a cross-species equivalence or DFU transfer (Tables~\ref{tab:p2mouse} and~\ref{tab:p2mousecounts}).

The paired baseline and predicted distances show the absolute error scale within each mouse arm (Figure~\ref{fig:p2f7}B); relative improvement alone would hide these differences.

\paperfigure{figure7_mouse_arms}{fig:p2f7}{Mouse POD7 computational audit for the source-labelled NDB, PDB and GDB arms. POD7 is excluded from fitting and scaling; rank-time fields trained at POD0, POD2 and POD30 predict from the observed POD2 source. (A) Relative energy-distance improvements over stay-still are −1.5\%, +32.1\% and +50.6\%. Hatching marks PDB, whose POD7 target fails the endpoint-distance ordering diagnostic. (B) Within each arm, the left grey bar is stay-still and the right coloured bar is prediction; hatching carries the same diagnostic meaning. The screen compares endpoint distances and does not prove or disprove geodesic interpolation. Each arm has one animal per time point and 4,182/7,002 mapped panel genes. Bars have no uncertainty intervals; cell counts do not supply independent animal replication or cross-species clinical validation.}


\FloatBarrier
\section{Discussion}
A fitted population field brought day-1 cells closer to the observed day-7 distribution in this acute human wound cohort. Improvement was present when day 7 was withheld globally and under a separate protocol that withheld each donor. These comparisons support distributional reconstruction for the measured interval in three healthy volunteers. Their distinction remains essential: transfer to a new donor can use day-7 observations from other people, whereas global timepoint holdout cannot. Neither task evaluates diabetic-foot-ulcer prognosis or the subsequent course of an individual chronic wound.

The mathematical interpretation is narrower than a recovered biological trajectory. Training samples independent cells from the endpoints of a donor's interval and regresses their linear-path velocities. It imposes neither a cell-lineage match nor a minimum-cost coupling. At the ideal regression solution, probability conservation concerns that specified path and sampling measure. It does not guarantee the distribution at an unobserved biological time, and empirical endpoints alone do not supply all the regularity needed for an exact characteristic-flow interpretation. This distinction is consistent with the fundamental ambiguity of dynamics inferred from snapshots and with the stronger assumptions required by identification theory \citep{weinreb2018,lavenant2021}.

The time-axis comparison exposes a consequential modeling choice. Linear days performed poorly at the target despite favorable results under several alternatives. Changing the axis here changes the straight bridges, their velocity scales and the optimization problem; the experiment is not a coordinate transformation of one already fitted field. Compression of the earliest interval is a possible contributor, but its separate effect was not isolated. The comparison therefore supports sensitivity to the chosen clock, not a preferred biological clock. Selecting an axis after inspecting the same held-out target also requires independent evaluation before claiming predictive validation \citep{cawley2010}.

Conditioning on a donor's source distribution improved training-donor fit but did not improve the overall held-out mean relative to the shared field. This separates additional fitting capacity from transfer performance. A shared field can give different velocities at different latent positions, so differences between donor-average displacements do not by themselves prove that donor context is necessary. The non-neural predictors also improved on unchanged source, indicating that part of the available structure can be captured by simple displacement rules. Broader comparisons of context-dependent and transport-based predictors should retain these simple controls \citep{bunne2023}.

The training and evaluation measures deserve explicit attention. Donor--segment cycling approximately balances donors during fitting, whereas the original pooled error weights their cells by yield. Re-evaluation of the same fixed prediction with equal donor mass preserved the improvement, and each donor separately improved under that pooled time-holdout field. The pooled result is therefore not explained solely by the observed cell-count imbalance. This sensitivity calculation does not refit a donor-balanced scaler, measure performance in a new donor, or remove dependence on the transferred expression coordinates.

Numerical integration and evaluation sampling were comparatively stable for the tested fixed field. Increasing RK4 resolution produced negligible changes, and repeated cell subsampling preserved the improvement. These checks reduce concern about those specific computational choices while leaving representation and training uncertainty open. Individual training configurations varied substantially across seeds. The aggregate gain from adding a second training donor persisted, but two training sizes cannot identify a learning-curve plateau. Overlapping training sets also correlate held-out errors, limiting the interpretation of a bootstrap interval based on only three donor blocks \citep{bengio2004}.

The principal limit is biological replication. Four times from three healthy donors describe a restricted acute-wound setting, and many cells do not increase the number of independent people \citep{zimmerman2021}. Each mouse arm has one animal per time point, so differences between times also include differences between animals. Incomplete gene coverage and variable performance further restrict cross-species interpretation. The next biological test should expand independent-donor evaluation under a frozen representation, clock and evaluation rule. An explicitly longitudinal diabetic-foot-ulcer cohort with authoritative outcomes would then be needed to test clinical staging or prognosis.

\section{Methods}
\subsection{Data, representation and join integrity}
GSE241132 contains three healthy acute-wound donors, labelled PWH26, PWH27 and PWH28, with Skin, Wound1, Wound7 and Wound30 samples \citep{liu2025_wound,geo_gse241132}. Among 65,521 loaded cells, 58,823 joined the author metadata, and 12,259 author-annotated fibroblasts entered the temporal analysis. The fibroblast counts were 2,152, 1,012, 3,821 and 5,274 at the four times, respectively. Annotation is joined using sample identity and original cell-row identity before quality filtering; a filtered positional row number is not a valid substitute. Earlier exploratory outputs based on positional mismatching are invalid and contribute no results here.

The expression encoder is frozen from GSE165816 \citep{theocharidis2022}; the human wound cohort was not used to refit it. It maps panel-normalized counts into a $q=15$ dimensional Gaussian-encoder mean $\mu$. Only 5,383 of 7,002 panel genes are covered in the human wound input (76.9\%); absent panel genes are zero-filled. The flow acts on $\mu\in\R^{15}$, not on simplex topic weights $\theta=\operatorname{softmax}(\mu)$. Successful reconstruction is therefore conditional on this incomplete transferred panel and this coordinate system.

For training-cell set $\mathcal T$, scaling parameters are fitted coordinatewise on $\mathcal T$ only:
\begin{equation}
a_j=\frac1{|\mathcal T|}\sum_{i\in\mathcal T}\mu_{ij},\qquad
b_j=\left[|\mathcal T|^{-1}\sum_{i\in\mathcal T}(\mu_{ij}-a_j)^2\right]^{1/2},\qquad
z_{ij}=\frac{\mu_{ij}-a_j}{\widetilde b_j}.\label{eq:p2scale}
\end{equation}
In Equation~\eqref{eq:p2scale}, $\widetilde b_j=b_j$ for $b_j>0$ and $\widetilde b_j=1$ if $b_j=0$. The identical transformation is applied to source and target evaluation cells. The target is excluded from estimating $a_j,b_j$, and all cells of the held-out donor are excluded in donor transfer. Distances from different training folds are consequently expressed in their respective fold-standardized spaces.

\subsection{Two distinct held-out tasks}
\textbf{Timepoint reconstruction.} All Wound7 cells are withheld from fitting and scaling. Training segments are Skin$\rightarrow$Wound1 and Wound1$\rightarrow$Wound30 within each donor. The fitted field advances the observed Wound1 source to the specified Wound7 coordinate. The pooled analysis uses three donors. Three additional runs omit one donor and evaluate Wound7 in the remaining donors; these are sensitivity analyses, not predictions for the omitted donor.

\textbf{Donor transfer.} A donor is excluded entirely from fitting and scaling; all four time points in the other two donors are retained. Adjacent within-donor segments Skin$\rightarrow$Wound1, Wound1$\rightarrow$Wound7 and Wound7$\rightarrow$Wound30 are used. The held donor's Wound1 cells initialize prediction, while their Wound7 cells are used for evaluation only. Because Wound7 is observed in training donors, this task does not test reconstruction of a globally unobserved time point. The method benchmark and donor-count experiment use this latter design. Table~\ref{tab:p2protocols} lists the exclusions and settings.

\subsection{Time coordinates and conditional flow matching}
Let $P_t^{(d)}$ denote the empirical distribution of standardized fibroblast coordinates for donor $d$ at day $t$. Write $s=g(t)$ for transformed time. Rank time assigns $(0,1/3,2/3,1)$ to days $(0,1,7,30)$, and the other tested transforms are
\begin{equation}
g_{\rm days}(t)=t/30,\qquad
g_{\rm sqrt}(t)=\sqrt{t/30},\qquad
g_{\rm log}(t)=\log(1+t)/\log31.\label{eq:p2times}
\end{equation}
Equation~\eqref{eq:p2times} defines the three day-based alternatives to rank time. For one training donor and segment $(s_a,s_b)$, endpoints are sampled \emph{independently} from the source and target cells of that donor. Thus the training coupling is the empirical product $\pi_{d,ab}=P_{t_a}^{(d)}\otimes P_{t_b}^{(d)}$. It is not an optimal-transport coupling and does not pair the same biological cell across times. With $\lambda\sim U(0,1)$,
\begin{equation}
z_\lambda=(1-\lambda)z_a+\lambda z_b,\quad
s_\lambda=(1-\lambda)s_a+\lambda s_b,\quad
u_{ab}=\frac{z_b-z_a}{s_b-s_a}.\label{eq:p2bridge}
\end{equation}
Equation~\eqref{eq:p2bridge} defines the interpolated state, time and target velocity. The implemented loss cycles over donor--segment pairs and averages squared error over cells and coordinates:
\begin{equation}
\mathcal L_{\rm CFM}(\phi)=\E_{(d,a,b),\pi_{d,ab},\lambda}
\left[\frac1q\left\|v_\phi(z_\lambda,s_\lambda)-u_{ab}\right\|_2^2\right].\label{eq:p2cfm}
\end{equation}
In Equation~\eqref{eq:p2cfm}, the division of the target velocity by $s_b-s_a$ is necessary: the target is velocity per unit transformed time. In particular, compressing the Skin-to-Wound1 segment on the linear-day axis increases its numerical velocity scale. Equations~\eqref{eq:p2bridge}--\eqref{eq:p2cfm} implement conditional flow matching \citep{lipman2023,tong2023,lipman2024guide}. They do not solve the Benamou--Brenier minimum-kinetic-energy problem \citep{benamou2000} or establish globally optimal transport.

Following the neural ordinary-differential-equation formulation \citep{chen2018ode}, predictions solve an initial-value problem and push forward the empirical source measure:
\begin{equation}
\frac{dz}{ds}=v_\phi(z,s),\qquad z(s_a)=z_a,\qquad
\widehat P_{s_b}=(\Phi^\phi_{s_a\rightarrow s_b})_\#\widehat P_{s_a}.\label{eq:p2ode}
\end{equation}
Equation~\eqref{eq:p2ode} defines prediction from the observed source distribution. The network uses two 128-unit SiLU hidden layers and a 32-dimensional sinusoidal time embedding. Training uses Adam with learning rate $10^{-3}$, 8,000 steps and batches of 256. Direct target predictions use 50 fixed fourth-order Runge--Kutta steps. The time-axis diagnostic instead uses 60 steps over the full Wound1-to-Wound30 segment and reads the grid point nearest the transformed Wound7 coordinate; it does not select the grid point with the lowest target distance. A separate fixed-field resolution comparison is described below; finite-step agreement does not prove convergence to an exact solution.

\subsection{What distributional success does and does not identify}
At the ideal population minimizer of the squared loss, the field is a conditional expectation of path velocities, $v^*(z,s)=\E[u_{ab}\mid z_\lambda=z,s_\lambda=s]$, where defined. Under suitable integrability and regularity, the induced marginal path satisfies a continuity equation,
\begin{equation}
\partial_s p_s+\nabla\!\cdot(p_sv^*)=0.\label{eq:p2continuity}
\end{equation}
Equation~\eqref{eq:p2continuity} describes conservation of probability mass in the model. It neither tracks proliferation/death rates nor determines individual-cell ancestry. Multiple couplings and fields can interpolate the same marginals; finite neural optimization adds approximation error. Matching the observed target distribution therefore does not identify a unique biological vector field \citep{lavenant2021}. With only four measurement times, the values of $g$ between observations are analysis choices rather than an identified biological clock.

\subsection{Distance estimand, finite-sample statistic and thresholds}
For independent $X,X'\sim P$ and $Y,Y'\sim Q$ with finite first moments, define energy distance without the optional square root \citep{szekely2013}:
\begin{equation}
D_E(P,Q)=2\E\|X-Y\|_2-\E\|X-X'\|_2-\E\|Y-Y'\|_2.\label{eq:p2ed}
\end{equation}
In Euclidean space, Equation~\eqref{eq:p2ed} is nonnegative and zero exactly when $P=Q$. The implementation uses the V-statistic, including diagonal zero distances:
\begin{align}
\widehat D_E(P_m,Q_n)={}&\frac{2}{mn}\sum_{i=1}^{m}\sum_{j=1}^{n}\|x_i-y_j\|_2
-\frac1{m^2}\sum_{i=1}^{m}\sum_{i'=1}^{m}\|x_i-x_{i'}\|_2\notag\\[-2pt]
&-\frac1{n^2}\sum_{j=1}^{n}\sum_{j'=1}^{n}\|y_j-y_{j'}\|_2.\label{eq:p2edv}
\end{align}
Equation~\eqref{eq:p2edv} is the exact distance between the two empirical measures, not an unbiased U-statistic for the population distance. For baseline $B=\widehat D_E(P_{\rm source},P_{\rm target})>0$ and predicted distance $D$, gain and relative improvement are
\begin{equation}
G=B-D,\qquad I=G/B=1-D/B.\label{eq:p2improvement}
\end{equation}
Equation~\eqref{eq:p2improvement} defines the absolute gain and the fraction reported as a percentage. Stay-still means $P_{\rm prediction}=P_{\rm source}$. Seeded subsampling caps the quadratic distance calculation (1,500 cells per distribution in the initial temporal and donor-transfer runs; the common benchmark uses its recorded 1,200-cell cap). Counts, caps and per-fold values accompany each comparison; independently run analyses need not have identical stochastic point estimates.

The target is split once into seeded disjoint halves, defining $N=\widehat D_E(P_{\rm target,1},P_{\rm target,2})$. The initial timepoint report asks only whether $G>0$, while the time-axis, donor-transfer, method and donor-count comparisons use $G>N$. The latter is a sampling-scale heuristic, not a significance test with known type-I error, a repeated-resampling interval, or an error bound. Selecting the best position on a path using target distances is retrospective; the preassigned Wound7 coordinate supplies the held-out result, and the full scan is a diagnostic rather than a tuned clinical predictor.

\subsection{Donor context and the non-neural predictors}
For a source sample, context concatenates the coordinatewise mean and standard deviation, $c_d=(\bar z_d,\operatorname{sd}(z_d))\in\R^{2q}$. The conditioned field $v_\phi(z,s,c_d)$ uses this context from the training segment's source. At test time, the held donor's source alone supplies $c_d$; its target is not used. Context is zeroed in 20\% of training batches. Both neural fields use the same donor-segment schedule and endpoint-sampling rule, although successive pseudorandom draws need not be identical. Six training-donor evaluations are reported separately from the three held-out-donor folds.

Let $\delta_d=\bar z_{d,7}-\bar z_{d,1}$ for a training donor. Mean displacement predicts $z+\overline\delta$; nearest-donor displacement predicts $z+\delta_{d^*}$, where $d^*$ maximizes cosine similarity between the source contexts. These are translations, not fitted neural fields. The entropic transport baseline \citep{cuturi2013} uses pooled training-source and training-target cells, uniform empirical marginals $a,b$, and normalized squared Euclidean cost $C$:
\begin{equation}
\Gamma_\varepsilon=\arg\min_{\Gamma\geq0,\,\Gamma\mathbf1=a,\,\Gamma^{\mathsf T}\mathbf1=b}
\left\{\langle C,\Gamma\rangle+\varepsilon\sum_{ij}\Gamma_{ij}(\log\Gamma_{ij}-1)\right\},\quad\varepsilon=0.05.\label{eq:p2ot}
\end{equation}
Equation~\eqref{eq:p2ot} defines the entropic transport plan fitted to the training endpoints. Sinkhorn iteration is capped at 2,000 steps, with at most 1,500 cells per endpoint. Row-normalized barycentres define training-source displacements; each held-out source cell receives the average displacement of its 15 nearest training-source neighbours. This is an out-of-sample displacement extension of a transport plan, not transport fitted to the held target.

\subsection{Donor-count estimand and dependence-aware summaries}
For each held donor $d$ and seed $r\in\{0,1,2\}$, let $D_{1,d,r}$ average the errors from the two possible single-training-donor choices, and let $D_{2,d,r}$ be the error with both available training donors. Equal weighting then gives
\begin{equation}
R_k=\frac1{3\cdot3}\sum_{d=1}^{3}\sum_{r=0}^{2}D_{k,d,r},\qquad
\Delta_d=\frac13\sum_{r=0}^{2}(D_{1,d,r}-D_{2,d,r}),\qquad
\overline\Delta=\frac13\sum_d\Delta_d.\label{eq:p2curve}
\end{equation}
Equation~\eqref{eq:p2curve} defines the seed-averaged estimand and donor-block differences. The plotted main curve uses seed 0, with other seeds shown separately. The 4,000-draw percentile interval resamples the three $\Delta_d$ blocks after averaging seeds \citep{efron1993bootstrap}. Overlapping training donors still induce dependence between held-out errors, so even this interval is a descriptive three-block summary, not a population coverage guarantee. Seeds and alternative training subsets are not additional biological replicates. Configuration-specific ranges quantify stochastic instability. Two values of $k$ cannot identify a plateau, a functional learning-curve model or a required number of donors.

\subsection{Geometry and mouse computation}
Donor displacement alignment is $\cos(\delta_d,\delta_e)=\delta_d^{\mathsf T}\delta_e/(\|\delta_d\|\|\delta_e\|)$. The time-versus-donor contrast compares 18 within-donor temporal distances with 12 between-donor same-time distances. Pair-distance bootstrap intervals are descriptive because pairs share donors. The geometric screen used for endpoints $a,b$ and held time $h$ is
\begin{equation}
J=\ind\{\widehat D_E(P_a,P_b)\geq\max[\widehat D_E(P_a,P_h),\widehat D_E(P_h,P_b)]\}.\label{eq:p2geometry}
\end{equation}
Equation~\eqref{eq:p2geometry} is a distance-ordering diagnostic. It does not imply equality in a triangle inequality, identify a geodesic, or establish that a target lies on a unique direct path. We report it as such rather than as proof of interpolability.

GSE326622 is evaluated separately for the source-labelled NDB, PDB and GDB mouse arms \citep{geo_gse326622}. POD7 is excluded from fitting and scaling; POD0, POD2 and POD30 provide the training time points, with POD2 supplying the prediction source. Each mouse fit uses rank time, 6,000 optimization steps and batches of 256. The human panel's mouse-symbol mapping covers 4,182 of 7,002 genes (59.7\%) in each arm, and each arm has one animal per time point. Animal identity and time are therefore confounded. The three arm computations assess protocol behaviour and representation transfer; cell numbers do not supply independent animals or cross-species clinical validation.

\subsection{Numerical-resolution and evaluation-sampling sensitivity}
The saved human expression coordinates and annotations supply a separate seed-0 pooled fit using the stated training exclusions, 8,000 optimization steps and batch size 256. The fitted field is then fixed. Predictions integrate the same 1,012 day-1 cells with 25, 50, 100 and 200 RK4 steps to the prespecified rank-time day-7 coordinate. Distance calculations use one fixed set of 1,500 day-7 cells and double-precision pairwise distances. Root-mean-square coordinate discrepancies are measured relative to the 200-step result (Table~\ref{tab:p2numerics}).

Evaluation sensitivity fixes the 50-step prediction and varies only cell subsampling. At each count of 250, 500 and 1,000 cells per distribution, 50 seeded draws without replacement use common source and target indices for predicted and unchanged-source distances. Each draw additionally samples two disjoint target subsets of that size for the reference distance. We report median improvement, the central 95\% range across draws and the count exceeding the reference (Table~\ref{tab:p2sampling}). These are conditional Monte Carlo summaries. They neither refit the field per draw nor estimate uncertainty across new donors.


\section*{Data and reproducibility}
Human acute-wound data and author cell annotations are available as GSE241132 \citep{liu2025_wound,geo_gse241132}. GSE165816 supplies the independently frozen expression model \citep{theocharidis2022}; its DFU clinical labels are not used to train the human temporal fields. Mouse data are from GSE326622 \citep{geo_gse326622}. This secondary computational study recruited no participants and performed no new animal experiments. Source-study consent and ethical approvals remain those of the original investigators. Sample identities, preprocessing exclusions, numerical settings and all reported comparisons are defined in the methods and supplementary tables.
Analysis code is released under the MIT License. Manuscripts, figures and source-study data remain under their respective rights and repository terms. The accompanying software archive, authored by Zeyu Fu, is available at \url{https://github.com/PeterPonyu/wound-population-dynamics}.
\FloatBarrier
\bibliographystyle{abbrvnat}
\bibliography{references}
\supplement
\section{Protocol and estimator specification}
Table~\ref{tab:p2protocols} separates the missing-time and missing-donor questions. Training endpoints are paired by donor, while the individual cells drawn from those endpoints are independent. Holding out a sample in an already trained donor is not equivalent to holding out that donor. The representation is frozen before all temporal comparisons.
\begin{table}[!htbp]\centering\small
\caption{Training and evaluation information for the six temporal protocols. Human data comprise three donors at Skin, Wound1, Wound7 and Wound30; mouse arms use postoperative days (POD) 0, 2, 7 and 30. Scaling follows the listed training exclusions. Source-only context at test time is permitted for donor-conditioned prediction, while the evaluation target is excluded from fitting in each held-out task. The explicit target column distinguishes prediction for an excluded donor from sensitivity analyses evaluated only in retained donors.}\label{tab:p2protocols}
\begin{tabularx}{\linewidth}{@{}p{36mm}LL@{}}\toprule
Protocol & Training data & Evaluation target and interpretation \\\midrule
Global Wound7 holdout & Skin, Wound1 and Wound30 from the included donors & Their Wound7 distributions; one globally missing measurement time. \\
Drop-one-donor sensitivity & Same time exclusions after removing one donor & Wound7 in the retained donors; sensitivity to donor composition, not prediction of the omitted donor. \\
Donor transfer & All four times in the other two donors & Omitted donor's Wound7 from their Wound1 source; observed interval in the training cohort. \\
Method benchmark & Same donor-transfer folds & Shared/conditioned fields and three non-neural predictors; training-donor fits shown separately. \\
Donor-count curve & One or two of the available training donors & Same held donor and target; seeds and subsets do not add independent donors. \\
Mouse POD7 audit & POD0, POD2 and POD30 in one source arm & POD7 from POD2; one animal per time point confounds time and individual. \\\bottomrule
\end{tabularx}\end{table}

\section{Numerical prediction sequence}
The reproducible evaluation order is: (1) resolve cell/sample identities and select the lineage; (2) define the held-out cells; (3) fit standardization only on training cells; (4) sample training endpoints and optimize the specified objective; (5) transport source cells with the fixed numerical integrator; (6) compute target distance, stay-still distance and the single split-half reference; (7) aggregate within the stated donor, seed and subset units. Target-distance path scans are diagnostics after prediction and are not used to select the reported target coordinate.

For completeness, the fixed-step RK4 update for step $h$ is
\begin{align}
k_1&=v(z_n,s_n),& k_2&=v(z_n+hk_1/2,s_n+h/2),\notag\\
k_3&=v(z_n+hk_2/2,s_n+h/2),& k_4&=v(z_n+hk_3,s_n+h),\notag\\
z_{n+1}&=z_n+\frac h6(k_1+2k_2+2k_3+k_4),&h&=(s_b-s_a)/50.\label{eq:p2rk4}
\end{align}
Equation~\eqref{eq:p2rk4} gives the numerical update for the direct target predictions. The conditioned field holds the source-derived context fixed throughout this integration. Under standard smoothness assumptions RK4 has fourth-order global convergence, but this property does not establish the actual error of the fitted field or certify the chosen step count in these experiments.

\section{Limits of the noise and geometric checks}
The finite empirical energy distance is nonzero when two independently sampled subsets represent the same population. The split-half reference therefore indicates one realized scale of target-sampling variation. It changes with cell number, latent scaling and random split, and does not include model-fit uncertainty. Comparing a prediction gain to it is a practical diagnostic with no established frequentist size in this analysis.

Likewise, the endpoint-distance ordering condition cannot show that a point lies on a line or a geodesic. In a metric space, even being closer to both endpoints than the endpoints are to each other is only a necessary-looking geometric screen, not an interpolation identity; here the unsquared-root energy quantity is used by convention and no triangle-inequality argument is invoked. PDB is consequently labelled as failing the ordering screen, not as a proven off-path biological state.

The fitted transport moves probability mass between empirical distributions; it does not model absolute cell expansion, death, migration into a biopsy or clinical closure. These unobserved processes and the non-uniqueness of couplings prevent a mechanistic interpretation of an individual streamline.

\section{Uncertainty and numerical results}
Tables~\ref{tab:p2folds}--\ref{tab:p2geometrytable} provide target errors, time-axis comparisons, the six-entry benchmark including stay-still, donor-specific method errors, the seed-averaged donor curve, mouse-arm contrasts and descriptive geometry intervals. All five predictive methods meeting a criterion in three folds is not five independent replications, because they share test donors. The lowest observed mean is reported without a population-level method ranking.

In the donor curve, two alternative single-donor fits contribute to each $k=1$ donor/seed value. The three-seed mean improvement is then aggregated over three held donors. Because training sets overlap between folds, the block interval does not estimate calibrated population or predictive uncertainty. Extending the curve requires new independent donors before any plateau rule can be evaluated. There is no statistically identified minimum cohort size from the two observed $k$ values.
Table~\ref{tab:p2insample} reports all six training-donor comparisons. Tables~\ref{tab:p2seedcurve} and~\ref{tab:p2seedconfigs} separate mean-displacement controls, seed-level means and configuration-specific training variability. Table~\ref{tab:p2mousecounts} lists the cell counts behind each mouse computation; each cell group still represents only one animal. The retained-donor sensitivity rows in Table~\ref{tab:p2folds} evaluate the two donors left in the fit, not the excluded donor.
{\small
\begin{longtable}{@{}p{50mm}p{28mm}p{28mm}p{28mm}p{31mm}@{}}
\caption{Human Wound7 evaluation in GSE241132. ED denotes the empirical energy-distance V-statistic in training-standardized latent space; lower values indicate closer predicted and observed distributions. The pooled row excludes Wound7 globally. Drop-one rows exclude the named donor and evaluate pooled Wound7 only in the two retained donors. Transfer rows instead evaluate the excluded donor after fitting all four time points of the other donors. Improvement is 100 times (stay-still ED minus predicted ED) divided by stay-still ED. Split-half ED uses one seeded target split and is not a confidence limit. Every row is a point estimate; these dependent fits do not create additional donors.}\label{tab:p2folds}\\
\toprule
Evaluation & Stay-still ED & Predicted ED & Split-half ED & Improvement \\
\midrule\endfirsthead
\multicolumn{5}{l}{\tablename\ \thetable\ (continued)}\\
\toprule
Evaluation & Stay-still ED & Predicted ED & Split-half ED & Improvement \\
\midrule\endhead
\bottomrule\endfoot
Pooled time holdout & 1.4069 & 0.7054 & 0.0043 & 49.86\% \\*
Drop PWH26; retained pair & 1.3357 & 0.7434 & 0.0043 & 44.34\% \\
Drop PWH27; retained pair & 1.2780 & 0.6258 & 0.0128 & 51.04\% \\
Drop PWH28; retained pair & 1.5460 & 0.7638 & 0.0080 & 50.59\% \\
PWH26 transfer & 1.8221 & 0.7110 & 0.0171 & 60.98\% \\
PWH27 transfer & 1.5355 & 0.2786 & 0.0096 & 81.85\% \\*
PWH28 transfer & 1.0704 & 0.1634 & 0.0418 & 84.73\% \\
\end{longtable}
}

{\small
\begin{longtable}{@{}p{41mm}p{32mm}p{31mm}p{32mm}p{29mm}@{}}
\caption{Historical single-seed four-clock diagnostic under global Wound7 holdout; not the revised paired exact-time comparison. The reference distance is 1.3879 and the split-half distance is 0.0061. ED denotes empirical energy distance in training-standardized latent space. Wound7 time is the global transformed coordinate; target distance uses the nearest point of the 60-step Wound1-to-Wound30 scan. Improvement is 100 times gain divided by stay-still distance. The three donors are reused, endpoint draws were not paired across clocks, and the path minimum is not substituted for target error. No interval is shown.}\label{tab:p2axes}\\
\toprule
Time coordinate & Wound7 time & Target ED & Improvement & Diverged \\
\midrule\endfirsthead
\multicolumn{5}{l}{\tablename\ \thetable\ (continued)}\\
\toprule
Time coordinate & Wound7 time & Target ED & Improvement & Diverged \\
\midrule\endhead
\bottomrule\endfoot
rank & 0.6667 & 0.6994 & 49.61\% & False \\*
days & 0.2333 & 1.6411 & -18.24\% & False \\
sqrt days & 0.4830 & 0.7695 & 44.56\% & False \\*
log days & 0.6055 & 0.6344 & 54.29\% & False \\
\end{longtable}
}

{\small
\begin{longtable}{@{}p{25mm}p{23mm}p{42mm}p{26mm}p{26mm}p{23mm}@{}}
\caption{Paired exact-time clock audit: four transforms, three seeds, 12 fields, three healthy donors. Each fit uses 8,000 steps, batch 256, a training-only scaler excluding all day-7 cells, and common source/target identities capped at 1,200 per donor. ED is energy distance in the common standardized encoder coordinates. Means equally weight donors and seeds. The seed range covers three equal-donor means, not a biological confidence interval. The centroid comparator uses the same donor\textquotesingle s day-1/day-30 means with the relevant time fraction. The final column is the maximum absolute donor/seed ED change from 50 to 200 RK4 steps, not proof of solver convergence. CPU checkpoint reload reproduces the predictions; no clock is selected using the held target.}\label{tab:p2pairedclock}\\
\toprule
Clock & Flow ED & Flow seed range & Centroid ED & Unchanged ED & Max refinement change \\
\midrule\endfirsthead
\multicolumn{6}{l}{\tablename\ \thetable\ (continued)}\\
\toprule
Clock & Flow ED & Flow seed range & Centroid ED & Unchanged ED & Max refinement change \\
\midrule\endhead
\bottomrule\endfoot
rank & 0.8104 & 0.8037 to 0.8229 & 0.6667 & 1.4746 & 4.98e-08 \\*
days & 2.1542 & 1.9740 to 2.3100 & 0.9775 & 1.4746 & 2.13e-06 \\
sqrt days & 0.8746 & 0.8588 to 0.8924 & 0.7476 & 1.4746 & 2.78e-08 \\*
log days & 0.7725 & 0.7400 to 0.8120 & 0.6654 & 1.4746 & 7.57e-08 \\
\end{longtable}
}

{\small
\begin{longtable}{@{}p{68mm}p{34mm}p{34mm}p{31mm}@{}}
\caption{Common Wound1-to-Wound7 benchmark across the three held-out human donors. ED is empirical energy distance in each training-standardized fold; reported means equally weight donors. Gain is stay-still ED minus method ED. Pass count records gains above the single split-half target reference, a descriptive criterion rather than a significance test. CFM denotes conditional flow matching and OT the entropic optimal-transport displacement extension. All methods share test donors; point estimates without confidence intervals do not establish a population ranking.}\label{tab:p2methods}\\
\toprule
Method & Mean ED & Mean gain & Pass count \\
\midrule\endfirsthead
\multicolumn{4}{l}{\tablename\ \thetable\ (continued)}\\
\toprule
Method & Mean ED & Mean gain & Pass count \\
\midrule\endhead
\bottomrule\endfoot
Stay-still & 1.4760 & 0.0000 & 0/3 \\*
Shared CFM & 0.3947 & 1.0813 & 3/3 \\
Mean displacement & 0.4186 & 1.0574 & 3/3 \\
Nearest donor & 0.5147 & 0.9613 & 3/3 \\
Conditioned CFM & 0.6176 & 0.8584 & 3/3 \\*
Entropic OT extension & 0.7284 & 0.7476 & 3/3 \\
\end{longtable}
}

{\small
\begin{longtable}{@{}p{68mm}p{34mm}p{34mm}p{31mm}@{}}
\caption{Empirical energy distances for each method and held-out human donor in the common benchmark. Each column uses the other two donors for fitting and scaling, the named donor's Wound1 cells as source and its Wound7 cells only for evaluation. Lower is better; stay-still retains the source distribution. CFM means conditional flow matching; OT is the entropic transport displacement extension. Entries are point estimates with a 1,200-cell distance cap. Training draws and this cap differ from separately run transfer experiments, so those point estimates need not coincide.}\label{tab:p2methodsbydonor}\\
\toprule
Method & PWH26 & PWH27 & PWH28 \\
\midrule\endfirsthead
\multicolumn{4}{l}{\tablename\ \thetable\ (continued)}\\
\toprule
Method & PWH26 & PWH27 & PWH28 \\
\midrule\endhead
\bottomrule\endfoot
Stay-still & 1.8221 & 1.5355 & 1.0704 \\*
Shared CFM & 0.7110 & 0.2881 & 0.1850 \\
Mean displacement & 0.5674 & 0.3496 & 0.3388 \\
Nearest donor & 0.6792 & 0.4318 & 0.4331 \\
Conditioned CFM & 1.4730 & 0.1990 & 0.1809 \\*
Entropic OT extension & 0.7201 & 0.6228 & 0.8423 \\
\end{longtable}
}

{\small
\begin{longtable}{@{}p{53mm}p{38mm}p{38mm}p{38mm}@{}}
\caption{Seed-averaged donor-count comparison after returning each prediction to the frozen encoder space and scoring in a common reference within held donor. The reference uses both non-held donors and never enters a single-donor fit. ED is empirical energy distance; k is the number of training donors. Single-donor choices are averaged within held donor and seed, followed by seeds 0, 1 and 2. Reduction is k=1 minus k=2. The donor-balanced reduction is 0.1483, with a descriptive three-block percentile 95\% range [0.0732, 0.2727]. Overlapping training sets prevent interpreting this range as a population coverage guarantee.}\label{tab:p2curve}\\
\toprule
Held donor & k=1 ED & k=2 ED & Reduction \\
\midrule\endfirsthead
\multicolumn{4}{l}{\tablename\ \thetable\ (continued)}\\
\toprule
Held donor & k=1 ED & k=2 ED & Reduction \\
\midrule\endhead
\bottomrule\endfoot
PWH26 & 1.0313 & 0.7586 & 0.2727 \\*
PWH27 & 0.3359 & 0.2368 & 0.0991 \\*
PWH28 & 0.2726 & 0.1994 & 0.0732 \\
\end{longtable}
}

{\small
\begin{longtable}{@{}p{21mm}p{27mm}p{27mm}p{34mm}p{27mm}p{29mm}@{}}
\caption{Fibroblast cell counts and target-sampling reference for each mouse arm. POD means postoperative day. POD0, POD2 and POD30 enter training; POD7 is the held-out target. Every arm/timepoint cell group comes from one animal, so these counts quantify distribution sampling and not biological replication. The final column is one seeded split-half POD7 energy distance in the training-standardized space. NDB, PDB and GDB are retained source-study arm labels. All arms share 4,182/7,002 mapped panel genes.}\label{tab:p2mousecounts}\\
\toprule
Arm & POD0 & POD2 & POD7 held out & POD30 & Split-half ED \\
\midrule\endfirsthead
\multicolumn{6}{l}{\tablename\ \thetable\ (continued)}\\
\toprule
Arm & POD0 & POD2 & POD7 held out & POD30 & Split-half ED \\
\midrule\endhead
\bottomrule\endfoot
NDB & 2818 & 8177 & 12914 & 550 & 0.00544 \\*
PDB & 1197 & 471 & 2500 & 813 & 0.01091 \\*
GDB & 223 & 2400 & 11590 & 1402 & 0.00617 \\
\end{longtable}
}

{\small
\begin{longtable}{@{}p{22mm}p{36mm}p{36mm}p{34mm}p{37mm}@{}}
\caption{Mouse POD7 audit for the source-labelled NDB, PDB and GDB arms. Each field uses rank time and POD0, POD2 and POD30 for fitting/scaling, then predicts from POD2 to excluded POD7. ED is empirical energy distance; improvement is 100 times (stay-still ED minus predicted ED) divided by stay-still ED. Each arm covers 4,182/7,002 human-panel genes and has one animal per time point. The ordering check asks whether endpoint distance is at least both endpoint-to-target distances; it is not proof of a geodesic. Values have no confidence intervals, and cell counts do not provide animal replication.}\label{tab:p2mouse}\\
\toprule
Arm & Stay-still ED & Predicted ED & Improvement & Ordering check \\
\midrule\endfirsthead
\multicolumn{5}{l}{\tablename\ \thetable\ (continued)}\\
\toprule
Arm & Stay-still ED & Predicted ED & Improvement & Ordering check \\
\midrule\endhead
\bottomrule\endfoot
NDB & 1.2341 & 1.2525 & -1.49\% & Yes \\*
PDB & 2.4345 & 1.6526 & 32.12\% & No \\*
GDB & 1.3532 & 0.6682 & 50.62\% & Yes \\
\end{longtable}
}

{\small
\begin{longtable}{@{}p{72mm}p{30mm}p{33mm}p{32mm}@{}}
\caption{Descriptive geometry for three healthy acute-wound donors. The first estimate divides the median of 18 within-donor between-time energy distances by the median of 12 between-donor same-time distances. The second is the mean of three pairwise cosines between donor Wound1-to-Wound7 mean-displacement vectors. Endpoints are percentile 95\% limits from 10,000 resamples of pairwise distances within class or of cosine pairs. Pairs share donors, so these intervals do not provide independent donor replication or a general population ordering.}\label{tab:p2geometrytable}\\
\toprule
Summary & Estimate & Lower endpoint & Upper endpoint \\
\midrule\endfirsthead
\multicolumn{4}{l}{\tablename\ \thetable\ (continued)}\\
\toprule
Summary & Estimate & Lower endpoint & Upper endpoint \\
\midrule\endhead
\bottomrule\endfoot
time to donor median ratio & 8.4570 & 6.0173 & 13.1114 \\*
mean displacement cosine & 0.9477 & 0.9432 & 0.9536 \\
\end{longtable}
}

{\small
\begin{longtable}{@{}p{32mm}p{47mm}p{29mm}p{27mm}p{30mm}@{}}
\caption{All six training-donor comparisons for shared and donor-conditioned flow matching. The first column identifies the donor excluded from that fit; the second is a retained training donor whose Wound1-to-Wound7 prediction is evaluated. ED is empirical energy distance in the fold-standardized coordinates. Both source and target of the evaluated donor were available during fitting, so these are fitting-capacity comparisons rather than donor-transfer tests. Conditioned errors are lower in all six rows. The repeated use of three donors prevents interpreting the rows as six independent biological replicates.}\label{tab:p2insample}\\
\toprule
Excluded donor & Evaluated training donor & Stay-still ED & Shared ED & Conditioned ED \\
\midrule\endfirsthead
\multicolumn{5}{l}{\tablename\ \thetable\ (continued)}\\
\toprule
Excluded donor & Evaluated training donor & Stay-still ED & Shared ED & Conditioned ED \\
\midrule\endhead
\bottomrule\endfoot
PWH26 & PWH27 & 1.6257 & 0.1277 & 0.0369 \\*
PWH26 & PWH28 & 1.0837 & 0.0930 & 0.0305 \\
PWH27 & PWH26 & 1.7375 & 0.2483 & 0.0511 \\
PWH27 & PWH28 & 1.0227 & 0.1248 & 0.0286 \\
PWH28 & PWH26 & 1.7873 & 0.3089 & 0.0584 \\*
PWH28 & PWH27 & 1.6191 & 0.1190 & 0.0306 \\
\end{longtable}
}

{\small
\begin{longtable}{@{}p{49mm}p{17mm}p{35mm}p{35mm}p{29mm}@{}}
\caption{Donor-count means in common reference coordinates. CFM denotes conditional flow matching; ED is empirical energy distance. The three CFM seeds reset both endpoint and network random generators. Deterministic controls are listed once. Full target-marginal prediction assigns equal mass to each selected training donor; matched prediction averages 50 draws with the same particle count as the held-source distribution. k=1 averages the two available single-training-donor choices before equally weighting held donors. Evaluation targets and scoring coordinates remain fixed across k and seed. Positive reductions favour k=2; these computational configurations retain three biological donors.}\label{tab:p2seedcurve}\\
\toprule
Predictor & Seed & k=1 mean ED & k=2 mean ED & Reduction \\
\midrule\endfirsthead
\multicolumn{5}{l}{\tablename\ \thetable\ (continued)}\\
\toprule
Predictor & Seed & k=1 mean ED & k=2 mean ED & Reduction \\
\midrule\endhead
\bottomrule\endfoot
Shared CFM & 0 & 0.5315 & 0.3854 & 0.1462 \\*
Mean displacement & Fixed & 0.5095 & 0.4186 & 0.0909 \\
Target marginal, full & Fixed & 0.2222 & 0.1676 & 0.0546 \\
Target marginal, matched & Fixed & 0.2332 & 0.1818 & 0.0514 \\
Unchanged source & Fixed & 1.4760 & 1.4760 & 0.0000 \\
Shared CFM & 1 & 0.5834 & 0.4276 & 0.1558 \\*
Shared CFM & 2 & 0.5250 & 0.3819 & 0.1430 \\
\end{longtable}
}

{\small
\begin{longtable}{@{}p{28mm}p{53mm}p{12mm}p{26mm}p{23mm}p{23mm}@{}}
\caption{Training stochasticity in common reference coordinates for every shared-field donor/training-subset configuration. Mean ED, sample SD and range summarize seeds 0, 1 and 2; range is maximum minus minimum. The observed source and target indices, target split and reference scaler are fixed within held donor, so variation within each row reflects fitting randomness. Seeds and overlapping training subsets are not independent donors. The single split-half target reference measures finite-cell sampling and omits fitting variation.}\label{tab:p2seedconfigs}\\
\toprule
Held donor & Training donors & k & Mean ED & SD & Range \\
\midrule\endfirsthead
\multicolumn{6}{l}{\tablename\ \thetable\ (continued)}\\
\toprule
Held donor & Training donors & k & Mean ED & SD & Range \\
\midrule\endhead
\bottomrule\endfoot
PWH26 & PWH27 & 1 & 1.2253 & 0.1920 & 0.3377 \\*
PWH26 & PWH28 & 1 & 0.8373 & 0.0056 & 0.0097 \\
PWH26 & PWH27|PWH28 & 2 & 0.7586 & 0.1238 & 0.2334 \\
PWH27 & PWH26 & 1 & 0.1952 & 0.0287 & 0.0573 \\
PWH27 & PWH28 & 1 & 0.4766 & 0.1098 & 0.1962 \\
PWH27 & PWH26|PWH28 & 2 & 0.2368 & 0.0516 & 0.1021 \\
PWH28 & PWH26 & 1 & 0.2468 & 0.0351 & 0.0655 \\
PWH28 & PWH27 & 1 & 0.2985 & 0.0816 & 0.1633 \\*
PWH28 & PWH26|PWH27 & 2 & 0.1994 & 0.0029 & 0.0057 \\
\end{longtable}
}

{\small
\begin{longtable}{@{}p{28mm}p{46mm}p{44mm}p{48mm}@{}}
\caption{Fixed-field Runge--Kutta resolution comparison for pooled day-7 holdout. A single seed-0 field was refitted with the specified 8,000 optimization steps and 256-cell batches on the saved expression coordinates, excluding day 7 from fitting and scaling. All rows use the same 1,012 day-1 cells and the same 1,500 sampled target cells. Energy distance is the empirical V-statistic evaluated in double precision; improvement is relative to the same unchanged-source baseline (1.406924). RMS is the root-mean-square coordinate difference from the 200-step solution; it is a numerical discrepancy, not a confidence interval. The comparison is conditional on one field and finite-precision arithmetic and does not establish formal convergence to an exact ODE solution.}\label{tab:p2numerics}\\
\toprule
RK4 steps & Energy distance & Improvement & RMS vs 200 steps \\
\midrule\endfirsthead
\multicolumn{4}{l}{\tablename\ \thetable\ (continued)}\\
\toprule
RK4 steps & Energy distance & Improvement & RMS vs 200 steps \\
\midrule\endhead
\bottomrule\endfoot
25 & 0.705395409 & 49.862585\% & 5.46e-07 \\*
50 & 0.705395449 & 49.862582\% & 4.76e-07 \\
100 & 0.705395463 & 49.862581\% & 5.36e-07 \\*
200 & 0.705395452 & 49.862582\% & 0 \\
\end{longtable}
}

{\small
\begin{longtable}{@{}p{34mm}p{17mm}p{37mm}p{42mm}p{37mm}@{}}
\caption{Evaluation-sampling sensitivity for the same fixed field with 50 RK4 steps. For each cell count, 50 seeded draws sample source and target cells without replacement; prediction and unchanged-source comparisons use identical source indices and identical target indices within a draw. An additional pair of disjoint target subsets at that count defines a split-half reference. Percentages summarize relative improvement over unchanged source; the central 95\% draw range measures Monte Carlo sensitivity conditional on this cohort and field, not patient-level confidence. The final column counts draws whose gain exceeds the split-half reference. Cell count and resampling do not increase the three-donor biological sample size.}\label{tab:p2sampling}\\
\toprule
Cells per distribution & Draws & Median improvement & Central 95\% draw range & Gain exceeds reference \\
\midrule\endfirsthead
\multicolumn{5}{l}{\tablename\ \thetable\ (continued)}\\
\toprule
Cells per distribution & Draws & Median improvement & Central 95\% draw range & Gain exceeds reference \\
\midrule\endhead
\bottomrule\endfoot
250 & 50 & 48.46\% & 44.24 to 51.57\% & 50/50 \\*
500 & 50 & 48.39\% & 46.22 to 51.69\% & 50/50 \\*
1000 & 50 & 49.26\% & 47.86 to 50.60\% & 50/50 \\
\end{longtable}
}

{\small
\begin{longtable}{@{}p{40mm}p{37mm}p{31mm}p{31mm}p{28mm}@{}}
\caption{Evaluation of the saved pooled 50-step field in training-standardized coordinates, using all 1,012 Wound1 and 3,821 Wound7 cells. ED is the double-precision weighted empirical energy distance, without subsampling. Cell-pooled weights are uniform; donor-balanced weights allocate one third to each donor and equal mass within donor. Predictions retain source weights. The final three rows evaluate training donors separately, not donor-transfer folds. Improvement uses the corresponding unchanged-source ED. These three-donor point estimates have no confidence intervals and differ from capped-distance estimates.}\label{tab:p2weights}\\
\toprule
Evaluation & Source / target cells & Stay-still ED & Predicted ED & Improvement \\
\midrule\endfirsthead
\multicolumn{5}{l}{\tablename\ \thetable\ (continued)}\\
\toprule
Evaluation & Source / target cells & Stay-still ED & Predicted ED & Improvement \\
\midrule\endhead
\bottomrule\endfoot
Cell-pooled & 1012 / 3821 & 1.36146 & 0.68890 & 49.40\% \\*
Equal donor mass & 1012 / 3821 & 1.31723 & 0.66265 & 49.69\% \\
PWH26 & 255 / 920 & 1.75803 & 0.92909 & 47.15\% \\
PWH27 & 497 / 1992 & 1.58457 & 0.84112 & 46.92\% \\*
PWH28 & 260 / 909 & 1.07483 & 0.64420 & 40.07\% \\
\end{longtable}
}

{\small
\begin{longtable}{@{}p{33mm}p{33mm}p{33mm}p{33mm}p{33mm}@{}}
\caption{Observed author-annotated fibroblasts retained after the sample--barcode join in GSE241132. Donors A, B and C correspond to source donor labels PWH26, PWH27 and PWH28. Each cell-count entry comes from one biopsy; the three donors supply the biological replication. The 12 entries sum to 12,259 cells and provide the material for expression and geometry displays. Day 0 is intact skin, followed by acute wounds at days 1, 7 and 30.}\label{tab:p2biologicalcounts}\\
\toprule
Donor & Day 0 cells & Day 1 cells & Day 7 cells & Day 30 cells \\
\midrule\endfirsthead
\multicolumn{5}{l}{\tablename\ \thetable\ (continued)}\\
\toprule
Donor & Day 0 cells & Day 1 cells & Day 7 cells & Day 30 cells \\
\midrule\endhead
\bottomrule\endfoot
Donor A & 724 & 255 & 920 & 843 \\*
Donor B & 521 & 497 & 1992 & 2096 \\*
Donor C & 907 & 260 & 909 & 2335 \\
\end{longtable}
}

{\small
\begin{longtable}{@{}p{17mm}p{42mm}p{24mm}p{24mm}p{24mm}p{32mm}@{}}
\caption{Retrospective global missing-time comparison in three donors. Each error is the equally weighted mean of three within-donor, full-cell empirical energy distances in all 15 training-standardized coordinates. The table gives its mean and range over seeds 0, 1 and 2; deterministic baselines have zero seed range. Improvement is relative to the corresponding unchanged-source error. Withheld day 7 uses day 1 as source and day 30 as anchor; withheld day 1 uses intact skin as source and day 7 as anchor. Both have rank-time interpolation fraction one half. The held-time cells are excluded from scaling, fitting and anchor construction. Seed ranges are not donor confidence intervals, and absolute errors from different tasks use different scalers.}\label{tab:p2interpolation}\\
\toprule
Held day & Predictor & Mean ED & Minimum ED & Maximum ED & Mean improvement \\
\midrule\endfirsthead
\multicolumn{6}{l}{\tablename\ \thetable\ (continued)}\\
\toprule
Held day & Predictor & Mean ED & Minimum ED & Maximum ED & Mean improvement \\
\midrule\endhead
\bottomrule\endfoot
7 & Unchanged source & 1.4725 & 1.4725 & 1.4725 & 0.00\% \\*
7 & Centroid translation & 0.6674 & 0.6674 & 0.6674 & 54.67\% \\
7 & Independent bridge & 0.9284 & 0.9224 & 0.9386 & 36.95\% \\
7 & Shared flow & 0.8110 & 0.8044 & 0.8237 & 44.93\% \\
1 & Unchanged source & 1.3460 & 1.3460 & 1.3460 & 0.00\% \\
1 & Centroid translation & 0.9293 & 0.9293 & 0.9293 & 30.96\% \\
1 & Independent bridge & 1.0959 & 1.0920 & 1.1007 & 18.58\% \\*
1 & Shared flow & 1.0571 & 0.9484 & 1.1231 & 21.46\% \\
\end{longtable}
}

{\small
\begin{longtable}{@{}p{17mm}p{24mm}p{48mm}p{24mm}p{24mm}p{24mm}@{}}
\caption{Every donor-specific error in the two global missing-time tasks, summarized over the three computational seeds. ED is the full 15-dimensional empirical energy distance using all available source-derived predictions and target cells in that donor. The mean, minimum and maximum describe seeds, not independent specimens. The centroid baseline has lower seed-averaged error than the field in all three day-7 donors and two of three day-1 donors. Each donor supplies both training endpoints in this retrospective design; these rows are not held-out-donor transfer.}\label{tab:p2interpolationdonors}\\
\toprule
Held day & Donor & Predictor & Mean ED & Minimum ED & Maximum ED \\
\midrule\endfirsthead
\multicolumn{6}{l}{\tablename\ \thetable\ (continued)}\\
\toprule
Held day & Donor & Predictor & Mean ED & Minimum ED & Maximum ED \\
\midrule\endhead
\bottomrule\endfoot
7 & Donor A & Unchanged source & 1.7580 & 1.7580 & 1.7580 \\*
7 & Donor A & Centroid translation & 0.7058 & 0.7058 & 0.7058 \\
7 & Donor A & Independent bridge & 0.8633 & 0.8508 & 0.8712 \\
7 & Donor A & Shared flow & 0.9069 & 0.8878 & 0.9291 \\
7 & Donor B & Unchanged source & 1.5846 & 1.5846 & 1.5846 \\
7 & Donor B & Centroid translation & 0.6898 & 0.6898 & 0.6898 \\
7 & Donor B & Independent bridge & 0.9808 & 0.9591 & 0.9980 \\
7 & Donor B & Shared flow & 0.8645 & 0.8411 & 0.8804 \\
7 & Donor C & Unchanged source & 1.0748 & 1.0748 & 1.0748 \\
7 & Donor C & Centroid translation & 0.6066 & 0.6066 & 0.6066 \\
7 & Donor C & Independent bridge & 0.9411 & 0.8981 & 0.9889 \\
7 & Donor C & Shared flow & 0.6614 & 0.6442 & 0.6868 \\
1 & Donor A & Unchanged source & 1.1341 & 1.1341 & 1.1341 \\
1 & Donor A & Centroid translation & 0.8704 & 0.8704 & 0.8704 \\
1 & Donor A & Independent bridge & 1.0311 & 1.0192 & 1.0471 \\
1 & Donor A & Shared flow & 0.8410 & 0.7370 & 0.8960 \\
1 & Donor B & Unchanged source & 1.5517 & 1.5517 & 1.5517 \\
1 & Donor B & Centroid translation & 1.0900 & 1.0900 & 1.0900 \\
1 & Donor B & Independent bridge & 1.3196 & 1.3132 & 1.3232 \\
1 & Donor B & Shared flow & 1.3462 & 1.2486 & 1.4146 \\
1 & Donor C & Unchanged source & 1.3522 & 1.3522 & 1.3522 \\
1 & Donor C & Centroid translation & 0.8275 & 0.8275 & 0.8275 \\
1 & Donor C & Independent bridge & 0.9369 & 0.9258 & 0.9525 \\*
1 & Donor C & Shared flow & 0.9841 & 0.8597 & 1.0585 \\
\end{longtable}
}

{\small
\begin{longtable}{@{}p{17mm}p{41mm}p{25mm}p{28mm}p{28mm}p{25mm}@{}}
\caption{Geometry controls for globally withheld days 7 and 1. Entries average seeds within each donor and then weight the three donors equally. ED is full-cell empirical energy distance in the 15 training-standardized coordinates. Centered ED subtracts each distribution\textquotesingle s own mean only for evaluation; it is a separate shape diagnostic and is not an additive error decomposition. Centroid error is Euclidean distance between means. Variance ratio is predicted total coordinate variance divided by observed target variance. Location--scale interpolates only the permitted endpoint means and standard deviations at rank fraction one half; no target mean or dispersion enters prediction. All entries are point estimates, not population uncertainty intervals.}\label{tab:p2shapediagnostics}\\
\toprule
Held day & Predictor & ED & Centered ED & Centroid error & Variance ratio \\
\midrule\endfirsthead
\multicolumn{6}{l}{\tablename\ \thetable\ (continued)}\\
\toprule
Held day & Predictor & ED & Centered ED & Centroid error & Variance ratio \\
\midrule\endhead
\bottomrule\endfoot
7 & Unchanged source & 1.4725 & 0.1495 & 2.5971 & 0.761 \\*
7 & Centroid translation & 0.6674 & 0.1495 & 1.5375 & 0.761 \\
7 & Independent bridge & 0.9284 & 0.2691 & 1.5426 & 0.319 \\
7 & Shared flow & 0.8110 & 0.1441 & 1.6835 & 0.543 \\
7 & Location-scale & 0.6982 & 0.1466 & 1.5375 & 0.622 \\
1 & Unchanged source & 1.3460 & 0.0647 & 2.5213 & 0.990 \\
1 & Centroid translation & 0.9293 & 0.0647 & 2.0270 & 0.990 \\
1 & Independent bridge & 1.0959 & 0.1445 & 2.0215 & 0.579 \\
1 & Shared flow & 1.0571 & 0.1220 & 2.0158 & 0.641 \\*
1 & Location-scale & 0.9034 & 0.0667 & 2.0270 & 1.107 \\
\end{longtable}
}

{\small
\begin{longtable}{@{}p{19mm}p{45mm}p{32mm}p{41mm}p{27mm}@{}}
\caption{Matched evaluation with 200 source and 200 target cells per donor, sampled without replacement in 50 draws with seed 29. Each draw shares indices across all predictors and fit seeds, then averages seeds within donor and donors equally. Entries are shared-flow error minus comparator error; positive differences favor the comparator. Full ranges describe observed-cell resampling sensitivity, not donor confidence intervals or hypothesis tests. The 50 draws still involve only the same three biological donors.}\label{tab:p2matchedsampling}\\
\toprule
Held day & Comparator & Mean difference & Full draw range & Comparator lower \\
\midrule\endfirsthead
\multicolumn{5}{l}{\tablename\ \thetable\ (continued)}\\
\toprule
Held day & Comparator & Mean difference & Full draw range & Comparator lower \\
\midrule\endhead
\bottomrule\endfoot
7 & Centroid translation & 0.1400 & 0.1132 to 0.1716 & 50/50 \\*
7 & Location-scale & 0.1102 & 0.0847 to 0.1405 & 50/50 \\
7 & Independent bridge & -0.1151 & -0.1576 to -0.0741 & 0/50 \\
1 & Centroid translation & 0.1243 & 0.0977 to 0.1545 & 50/50 \\
1 & Location-scale & 0.1497 & 0.1194 to 0.1763 & 50/50 \\*
1 & Independent bridge & -0.0407 & -0.0833 to -0.0070 & 0/50 \\
\end{longtable}
}

{\small
\begin{longtable}{@{}p{27mm}p{30mm}p{35mm}p{46mm}p{27mm}@{}}
\caption{Training-target marginal predictor compared with the historical shared-field benchmark. Each target-marginal draw samples the permitted training donors\textquotesingle{} Wound7 cells without replacement, matches the original held-source particle count (255, 497 or 260), and assigns each training donor exactly half of total probability mass. It uses no held-source expression. Scaling and target indices are identical to the original benchmark. Means and full ranges describe 50 draws; they are not donor confidence intervals. The shared-field column is its original saved error, because historical per-cell predictions were unavailable for a new paired evaluation. Across donors the matched marginal mean is 0.1818 versus 0.3947 for the field; lower is better.}\label{tab:p2targetmarginal}\\
\toprule
Held donor & Shared-field ED & Target-marginal ED & Full draw range & Marginal lower \\
\midrule\endfirsthead
\multicolumn{5}{l}{\tablename\ \thetable\ (continued)}\\
\toprule
Held donor & Shared-field ED & Target-marginal ED & Full draw range & Marginal lower \\
\midrule\endhead
\bottomrule\endfoot
PWH26 & 0.7110 & 0.1569 & 0.1147 to 0.2106 & 50/50 \\*
PWH27 & 0.2881 & 0.1416 & 0.1055 to 0.1769 & 50/50 \\*
PWH28 & 0.1850 & 0.2469 & 0.1860 to 0.3340 & 0/50 \\
\end{longtable}
}

{\small
\begin{longtable}{@{}p{31mm}p{27mm}p{29mm}p{33mm}p{35mm}@{}}
\caption{Retrospective source-replacement control for the nine saved two-training-donor fields (three held donors by three seeds). All scores are empirical energy distance in the same reference coordinates and on the same held-donor Wound7 target cells; lower is better. Real is the actual held-donor Wound1 distribution. Substituted sources are 50 without-replacement, source-count-matched draws from permitted training donors\textquotesingle{} Wound1 cells, integrated by the identical saved field, with one-half mass assigned to each training donor. The target marginal instead draws from their Wound7 cells without reading the held source; its draw IDs are not paired with substituted-source draws. Draws are averaged within field, then seeds within held donor; the final row equally weights three donors. The difference is substituted minus real, so a negative value favors the substituted input. These are retrospective, dependent computations on three healthy donors, not confidence intervals, new biological replicates or a model-selection test.}\label{tab:p2sourceswap}\\
\toprule
Held donor & Real-source ED & Substituted ED & Target marginal ED & Substituted minus real \\
\midrule\endfirsthead
\multicolumn{5}{l}{\tablename\ \thetable\ (continued)}\\
\toprule
Held donor & Real-source ED & Substituted ED & Target marginal ED & Substituted minus real \\
\midrule\endhead
\bottomrule\endfoot
PWH26 & 0.7586 & 0.2802 & 0.1569 & -0.4785 \\*
PWH27 & 0.2368 & 0.2732 & 0.1416 & 0.0363 \\
PWH28 & 0.1994 & 0.1816 & 0.2469 & -0.0179 \\*
Three-donor mean & 0.3983 & 0.2450 & 0.1818 & -0.1533 \\
\end{longtable}
}

{\small
\begin{longtable}{@{}p{35mm}p{33mm}p{33mm}p{33mm}p{33mm}@{}}
\caption{Exact target-marginal mixture decomposition in fixed held-fold coordinates. Separate ED is the mean of the two single-training-donor target errors; mixed ED scores their equal-mass empirical mixture. Reduction is separate minus mixed and equals one quarter of their mutual ED, independently of the held target. ED denotes the empirical energy discrepancy without its square root. Complete permitted training-Wound7 distributions and the original held-target indices reproduce the earlier full-marginal curve. Maximum identity residual is 0.0000000000000024. The final row equally weights three donors. This identity does not decompose a jointly fitted neural field, and lower mixture error does not itself identify target-relevant information gain.}\label{tab:p2mixtureidentity}\\
\toprule
Held donor & Separate ED & Mixed ED & Reduction & Mutual ED / 4 \\
\midrule\endfirsthead
\multicolumn{5}{l}{\tablename\ \thetable\ (continued)}\\
\toprule
Held donor & Separate ED & Mixed ED & Reduction & Mutual ED / 4 \\
\midrule\endhead
\bottomrule\endfoot
PWH26 & 0.202462 & 0.137777 & 0.064685 & 0.064685 \\*
PWH27 & 0.199092 & 0.135287 & 0.063805 & 0.063805 \\
PWH28 & 0.265045 & 0.229730 & 0.035315 & 0.035315 \\*
Three-donor mean & 0.222199 & 0.167598 & 0.054601 & 0.054601 \\
\end{longtable}
}

{\small
\begin{longtable}{@{}p{34mm}p{30mm}p{37mm}p{37mm}p{29mm}@{}}
\caption{Matched-particle source substitutions in nine existing two-training-donor fields (three held donors, seeds 0--2). Each score uses 200 predicted particles: real held-source cells, one individual permitted training-donor source, or a mixture of 100 cells from each permitted donor. Fifty without-replacement draws share fixed held-target indices and reuse source indices across model seeds. Wrong A/B list actual source donors; their identities vary across held folds, so the final columns average fold-specific substitutions, not a single common donor. Scores average draws within field, then seeds within held donor, then donors equally. The individual-substitution error average is 0.3075. These are point energy discrepancies in common reference coordinates, not confidence intervals or new biological replication. Mixture identities are separately checked on the actual 100-cell components, not by equating them to the 200-cell individual distributions. No model is refitted.}\label{tab:p2mixtureaudit}\\
\toprule
Held donor & Real-source ED & Wrong A ED & Wrong B ED & Mixed ED \\
\midrule\endfirsthead
\multicolumn{5}{l}{\tablename\ \thetable\ (continued)}\\
\toprule
Held donor & Real-source ED & Wrong A ED & Wrong B ED & Mixed ED \\
\midrule\endhead
\bottomrule\endfoot
PWH26 & 0.7578 & PWH27: 0.3003 & PWH28: 0.3745 & 0.2799 \\*
PWH27 & 0.2488 & PWH26: 0.3504 & PWH28: 0.3057 & 0.2877 \\
PWH28 & 0.2043 & PWH26: 0.2503 & PWH27: 0.2636 & 0.1904 \\*
Three-donor mean & 0.4036 & 0.3003 & 0.3146 & 0.2527 \\
\end{longtable}
}

{\small
\begin{longtable}{@{}p{26mm}p{9mm}p{53mm}p{24mm}p{25mm}p{28mm}@{}}
\caption{Observable-space errors against actual held-donor expression in the same three healthy volunteers. All 5,383 assayed frozen-panel genes are retained without a performance filter; 1,619 absent genes are not observed zeros. Each observed cell or decoded particle is normalized on the common panel before equal-cell averaging. Saved target uses the original held-target identities; all day 7 is a predeclared sensitivity. k counts permitted training donors, not test donors. Twenty-seven saved fits use three computational seeds; no model is trained. Averages weight seeds, then training subsets, then held donors equally. TV is total variation between mean profiles; JS is Jensen--Shannon divergence in natural-log units, not its square root; RMSE is the root mean square error over gene proportions. Lower is better. Observed and decoded baselines are distinct; training marginals use only permitted donors. These are mean-profile errors, not single-cell transitions, full distributional fidelity, absolute-count prediction or biological confidence intervals.}\label{tab:p2observable}\\
\toprule
Target set & k & Predictor & TV & JS (nats) & Gene RMSE \\
\midrule\endfirsthead
\multicolumn{6}{l}{\tablename\ \thetable\ (continued)}\\
\toprule
Target set & k & Predictor & TV & JS (nats) & Gene RMSE \\
\midrule\endhead
\bottomrule\endfoot
Saved target & 1 & Decoded shared flow & 0.2934 & 0.0746 & 0.000779 \\*
Saved target & 1 & Decoded persistence & 0.3529 & 0.0990 & 0.001128 \\
Saved target & 1 & Decoded mean displacement & 0.2882 & 0.0717 & 0.000759 \\
Saved target & 1 & Decoded training marginal & 0.2590 & 0.0636 & 0.000600 \\
Saved target & 1 & Observed persistence & 0.3421 & 0.0928 & 0.001268 \\
Saved target & 1 & Observed training marginal & 0.1279 & 0.0155 & 0.000370 \\
Saved target & 2 & Decoded shared flow & 0.2823 & 0.0709 & 0.000694 \\
Saved target & 2 & Decoded persistence & 0.3529 & 0.0990 & 0.001128 \\
Saved target & 2 & Decoded mean displacement & 0.2825 & 0.0703 & 0.000714 \\
Saved target & 2 & Decoded training marginal & 0.2546 & 0.0627 & 0.000571 \\
Saved target & 2 & Observed persistence & 0.3421 & 0.0928 & 0.001268 \\
Saved target & 2 & Observed training marginal & 0.1115 & 0.0120 & 0.000322 \\
All day 7 & 1 & Decoded shared flow & 0.2929 & 0.0744 & 0.000780 \\
All day 7 & 1 & Decoded persistence & 0.3522 & 0.0988 & 0.001126 \\
All day 7 & 1 & Decoded mean displacement & 0.2878 & 0.0716 & 0.000760 \\
All day 7 & 1 & Decoded training marginal & 0.2584 & 0.0634 & 0.000600 \\
All day 7 & 1 & Observed persistence & 0.3412 & 0.0923 & 0.001265 \\
All day 7 & 1 & Observed training marginal & 0.1270 & 0.0153 & 0.000369 \\
All day 7 & 2 & Decoded shared flow & 0.2818 & 0.0707 & 0.000695 \\
All day 7 & 2 & Decoded persistence & 0.3522 & 0.0988 & 0.001126 \\
All day 7 & 2 & Decoded mean displacement & 0.2820 & 0.0702 & 0.000714 \\
All day 7 & 2 & Decoded training marginal & 0.2541 & 0.0624 & 0.000571 \\
All day 7 & 2 & Observed persistence & 0.3412 & 0.0923 & 0.001265 \\*
All day 7 & 2 & Observed training marginal & 0.1103 & 0.0118 & 0.000320 \\
\end{longtable}
}

{\small
\begin{longtable}{@{}p{26mm}p{32mm}p{22mm}p{25mm}p{27mm}p{33mm}@{}}
\caption{Decoder reconstruction reference for the observed held-donor day-7 cells. The frozen encoder means of the target cells themselves are decoded, restricted to the same 5,383 observed panel genes and normalized per cell before averaging. Errors compare this reconstruction with actual mean observed expression, not with another decoded target. TV is total variation; JS is Jensen--Shannon divergence in nats. The missing-mass column reports mean decoder probability on the 1,619 unmeasured panel genes before renormalization. This reference uses held-target information: it is not a forecast, true expression, a competing predictor or a mathematical lower error bound. The two target sets overlap and do not add independent observations.}\label{tab:p2decoderreference}\\
\toprule
Held donor & Target set & Cells & TV & JS (nats) & Unmeasured mass \\
\midrule\endfirsthead
\multicolumn{6}{l}{\tablename\ \thetable\ (continued)}\\
\toprule
Held donor & Target set & Cells & TV & JS (nats) & Unmeasured mass \\
\midrule\endhead
\bottomrule\endfoot
PWH26 & Saved target & 920 & 0.2314 & 0.0553 & 1.59\% \\*
PWH26 & All day 7 & 920 & 0.2314 & 0.0553 & 1.59\% \\
PWH27 & Saved target & 1200 & 0.2549 & 0.0655 & 1.62\% \\
PWH27 & All day 7 & 1992 & 0.2537 & 0.0651 & 1.62\% \\
PWH28 & Saved target & 909 & 0.2426 & 0.0584 & 1.79\% \\*
PWH28 & All day 7 & 909 & 0.2426 & 0.0584 & 1.79\% \\
\end{longtable}
}


\end{document}
